\documentclass[a4paper,12pt]{article}
% =========================================================
% XeLaTeX + tiếng Việt
% =========================================================
\usepackage{fontspec}
\usepackage{polyglossia}
\setdefaultlanguage{vietnamese}
\setmainfont{Times New Roman}

% =========================================================
% Gói cơ bản
% =========================================================
\usepackage{amsmath,amssymb,amsthm}
\usepackage{geometry}
\geometry{left=2.5cm,right=2.5cm,top=2.5cm,bottom=2.5cm}
\usepackage{xcolor}
\usepackage{tikz}
\usepackage{pgfplots}
\usepackage{fancyhdr}
\usepackage{enumitem}
\usepackage[hidelinks]{hyperref}
\usepackage{multicol}
\usepackage{float}
\usepackage{etoolbox}
\usepackage{tcolorbox}
\usepackage{array}
\usepackage{booktabs}
\usepackage{longtable}
\usepackage{setspace}
\usepackage{titlesec}
\usepackage{fontawesome5}
\usepackage{caption}

\pgfplotsset{compat=1.18}
\usetikzlibrary{
    shadows, positioning, calc, shapes.geometric,
    decorations.pathreplacing, decorations.pathmorphing,
    arrows.meta, fadings
}

% =========================================================
% Màu sắc
% =========================================================
\definecolor{maincolor}{RGB}{0,102,204}
\definecolor{secondcolor}{RGB}{204,51,102}
\definecolor{lightbg}{RGB}{245,245,250}
\definecolor{theoremcolor}{RGB}{0,102,102}
\definecolor{defcolor}{RGB}{153,51,102}
\definecolor{questioncolor}{RGB}{204,102,0}
\definecolor{aiblue}{RGB}{60,90,180}
\definecolor{bizgreen}{RGB}{0,130,90}
\definecolor{hintviolet}{RGB}{120,70,170}
\definecolor{softorange}{RGB}{245,140,60}
\definecolor{dialoguecolor}{RGB}{80,80,80}

% =========================================================
% Định nghĩa phông chữ cho kịch bản điện ảnh
% =========================================================
\newfontfamily{\scriptfont}{Courier New}[
    Scale=0.95,
    Ligatures=TeX,
    WordSpace=1.2
]

% =========================================================
% Header / Footer
% =========================================================
\pagestyle{fancy}
\fancyhf{}
\fancyhead[L]{\textcolor{maincolor}{\textbf{Bài giảng: Thống kê mô tả và Thống kê suy luận}}}
\fancyhead[R]{\textcolor{maincolor}{\textbf{\thepage}}}
\fancyfoot[C]{\textit{Thống kê không phải để tính – Thống kê để quyết định}}

% =========================================================
% Định nghĩa môi trường
% =========================================================
\newtheorem{defn}{\color{defcolor}\bfseries Định nghĩa}[section]
\newtheorem{theorem}{\color{theoremcolor}\bfseries Định lý}[section]
\theoremstyle{remark}
\newtheorem{remark}{\bfseries Nhận xét}[section]

% =========================================================
% Box theo phong cách AI-first
% =========================================================
\newtcolorbox{aibox}{
    colback=blue!3,
    colframe=aiblue,
    title=\faRobot\ Góc nhìn AI-first,
    fonttitle=\bfseries,
    rounded corners,
    boxrule=0.8pt
}
\newtcolorbox{metabox}{
    colback=purple!4,
    colframe=hintviolet,
    title=\faLightbulb\ Ẩn dụ trực giác,
    fonttitle=\bfseries,
    rounded corners,
    boxrule=0.8pt
}
\newtcolorbox{bizbox}{
    colback=green!4,
    colframe=bizgreen,
    title=\faChartLine\ Bài toán kinh doanh,
    fonttitle=\bfseries,
    rounded corners,
    boxrule=0.8pt
}
\newtcolorbox{notebox}{
    colback=orange!5,
    colframe=softorange,
    title=\faExclamationCircle\ Lưu ý học thuật,
    fonttitle=\bfseries,
    rounded corners,
    boxrule=0.8pt
}

% =========================================================
% Tiêu đề bài giảng
% =========================================================
\title{
    \vspace{-1.5cm}
    {\color{maincolor}\textbf{\huge THỐNG KÊ MÔ TẢ VÀ THỐNG KÊ SUY LUẬN}}\\[0.2cm]
    {\Large \textit{Bài giảng đầu tiên – Phương pháp AI-first}}\\[0.3cm]
    {\large Từ mẫu đến tổng thể – Từ tính toán đến quyết định}
}
\author{
    \textcolor{maincolor}{\textbf{Giáo sư AI}}\\
    \small Phương pháp giáo dục mới trong kỷ nguyên công nghệ
}
\date{\today}

\begin{document}
\maketitle
\thispagestyle{fancy}
\tableofcontents

\newpage
\begin{center}
    {\Large\color{secondcolor}\textbf{Dẫn nhập}}\\[0.4cm]
    \textit{\textcolor{maincolor}{“Thống kê không phải để tính – Thống kê để quyết định.”}}
\end{center}
\vspace{0.5cm}

% =========================================================
% PHẦN MỞ ĐẦU: NGUYÊN VĂN HỘI THOẠI
% =========================================================
{\scriptfont
\noindent\textcolor{dialoguecolor}{\textbf{Nguyên văn cuộc hội thoại:}}

\vspace{0.3cm}

\noindent\textcolor{dialoguecolor}{Ok. Vậy thì cái buổi đầu tiên mà thầy trò làm việc thì chúng ta sẽ dự định là \textbf{viết cái quyển thống kê xác suất theo cái tinh thần là những cái gì nó AI nó làm tốt thì chúng ta không làm nữa} để cho dân toán chuyên nghiệp nó làm. Còn những cái gì mà chúng ta học xác suất thống kê á là chúng ta có thể hiểu cách tính các tham số thống kê, rồi tính xong rồi thì chúng ta hiểu cái cách mà dựa vào đó mà hành động, cái đó mới là cái quan trọng. Có nghĩa là đầu tiên nó phải trí tuệ nhân tạo nó làm: tính trung bình, kỳ vọng, phương sai, tính mẫu, dự đoán độ chính xác, độ tin cậy, kiểm định đúng sai tất tần tật. Chúng ta phải nắm được ý nghĩa, cách làm nhưng không mất thời gian để làm vì trí tuệ nhân tạo nó làm quá tốt rồi.}

\vspace{0.2cm}

\noindent\textcolor{dialoguecolor}{Nếu không phải là học bản chất thì chúng mày học toán ngu là vì không nắm được bản chất mà cứ đi làm những bài tập ra đáp số rồi thi nọ thi kia nhưng thực tế không biết gì về toán cả. Bây giờ nhờ có trí tuệ nhân tạo dữ liệu lớn, sư phạm tốt, chúng ta sẽ viết giáo trình hoàn toàn mới mà chưa bao giờ thấy trên internet hay của Tây, của ta viết.}

\vspace{0.2cm}

\noindent\textcolor{dialoguecolor}{Thầy lên lớp chỉ "chửi" mà không dạy gì cả nhưng lớp vẫn 78\% hoặc 50\% qua môn, trong khi các lớp khác dạy đủ phương pháp, phấn bảng, hội thảo nhưng tỷ lệ qua môn thấp hơn. Trong thời đại công nghệ, \textbf{ông thầy đóng vai trò thiết kế lộ trình học tập, chỉ cho sinh viên sử dụng AI để tự soạn bài, soạn slide, tự đặt câu hỏi}. Nếu không có người thầy hướng dẫn thì sinh viên không thể làm được, dù có nghĩ tới cũng không dùng được AI hiệu quả.}

\vspace{0.2cm}

\noindent\textcolor{dialoguecolor}{\textbf{Môn thống kê có hai phần là thống kê mô tả và thống kê suy luận}, hai cái này là một cặp, nếu không dẫn sang suy luận thì coi như chưa học thống kê. Thống kê mô tả là làm việc trên mẫu, đo đạc và tính toán các đặc trưng của mẫu. Thống kê là để dự đoán hoặc tính toán các tham số của một quần thể rất lớn, ví dụ như đếm cá dưới hồ hay cây rừng mà không thể đếm hết được. Chúng ta làm việc trên một nhóm nhỏ (mẫu) rồi từ đó suy luận ra nhóm lớn (tổng thể) bằng phương pháp thống kê suy diễn.}

\vspace{0.2cm}

\noindent\textcolor{dialoguecolor}{Ví dụ không thể đo chiều cao của 100 triệu dân Việt Nam nên người ta đo ở vài vùng với vài trăm người để suy ra chiều cao cả nước; mấy trăm người đó gọi là mẫu. Việc tính trung bình chiều cao của họ gọi là đặc trưng của mẫu, từ đó suy luận ra tham số của tổng thể. Trong kinh doanh, người người ta cần đo đạc xem bao nhiêu phần trăm khách hàng thích sản phẩm để quyết định sản xuất tiếp hay dừng lại dựa trên việc suy luận thị trường từ tham số tổng thể.}

\vspace{0.2cm}

\noindent\textcolor{dialoguecolor}{Người học thống kê phải đọc được các tham số giống như giám đốc đọc bảng cân đối kế toán để ra quyết định. Thầy viết giáo trình theo tinh thần nói được, làm được, đọc được và hiểu đúng bản chất. Trong thống kê có nhiều kiến thức toán sâu sắc, nếu không hiểu bản chất thì chỉ biết tra bảng mà không ứng dụng được. Khi dùng AI, nó sẽ mổ xẻ, trả lời cặn kẽ và đưa ra những ẩn dụ phong phú, ví dụ như đưa kiến thức về món ăn cho dễ hiểu.}

\vspace{0.2cm}

\noindent\textcolor{dialoguecolor}{Con AI sẽ đi thẳng vào bản chất, giải thích tại sao có công thức đó và ý nghĩa của nó thay vì vòng vo diễn giải cách tính. Tuy nhiên, cần có phần mềm và sự hướng dẫn của thầy để chuyển đề cương thành bài giảng, bài tập và sách; nếu không có trình độ mà hỏi AI vài câu lặt vặt thì nó sẽ phát hiện ra và trả lời sai, dẫn đến kết quả hỏng bét.}
}

\vspace{1cm}

% =========================================================
% PHẦN 1: GIỚI THIỆU PHƯƠNG PHÁP AI-FIRST
% =========================================================
\section*{\color{maincolor}I. Phương pháp AI-first – Học thống kê theo cách hoàn toàn mới}

Từ cuộc hội thoại trên, chúng ta rút ra được một triết lý giáo dục hoàn toàn mới: \textbf{những gì AI làm tốt, chúng ta không làm nữa}.

\begin{aibox}
\textbf{Tinh thần của phương pháp AI-first:}
\begin{itemize}
    \item AI tính toán: trung bình, kỳ vọng, phương sai, khoảng tin cậy, kiểm định... – AI làm quá tốt rồi.
    \item Con người hiểu \textbf{bản chất} và \textbf{ra quyết định}: đó mới là cái quan trọng.
    \item Người thầy đóng vai trò \textbf{thiết kế lộ trình học tập}, hướng dẫn sử dụng AI.
    \item Sinh viên được trang bị khả năng \textbf{đọc hiểu tham số} và \textbf{ứng dụng vào thực tế}.
\end{itemize}
\end{aibox}

\begin{center}
\begin{tikzpicture}[node distance=1cm]
    \node[draw, rounded corners, fill=maincolor!10, text width=3cm, align=center] (old) {Cách học cũ\\Tính toán thủ công};
    \node[draw, rounded corners, fill=red!20, text width=3cm, align=center, right=of old] (new) {Cách học mới\\Hiểu bản chất};
    \node[draw, rounded corners, fill=blue!20, text width=3cm, align=center, below=of new] (ai) {AI làm thay\\tính toán};
    
    \draw[-{Stealth[length=3mm]}, thick] (old) -- (new);
    \draw[-{Stealth[length=3mm]}, thick] (new) -- (ai);
\end{tikzpicture}
\captionof{figure}{Sự chuyển dịch trong phương pháp học thống kê thời AI.}
\end{center}

% =========================================================
% PHẦN 2: THỐNG KÊ MÔ TẢ
% =========================================================
\section*{\color{maincolor}II. Thống kê mô tả – Làm việc trên mẫu}

\begin{defn}
\textbf{Thống kê mô tả} là tập hợp các phương pháp dùng để \textbf{thu thập, tổ chức, trình bày và tính toán các đặc trưng} của dữ liệu mẫu.
\end{defn}

\begin{metabox}
Hãy tưởng tượng bạn là một \textbf{giám đốc kinh doanh}. Bạn muốn biết khách hàng của mình có hài lòng không. Bạn không thể hỏi tất cả hàng triệu khách hàng. Bạn chọn ra \textbf{500 khách hàng ngẫu nhiên} để khảo sát. 500 người đó gọi là \textbf{mẫu}.

\textbf{Thống kê mô tả} sẽ giúp bạn:
\begin{itemize}
    \item Tính \textbf{tỷ lệ hài lòng} trung bình của mẫu.
    \item Tính \textbf{độ phân tán} (mức độ chênh lệch giữa các câu trả lời).
    \item Vẽ \textbf{biểu đồ} để thấy xu hướng.
\end{itemize}
\end{metabox}

\subsection*{2.1 Các đặc trưng mẫu thường gặp}

\begin{table}[H]
\centering
\caption{Các đặc trưng thống kê mô tả cơ bản}
\begin{tabular}{@{}p{4cm} p{5cm} p{3cm}@{}}
\toprule
\textbf{Đặc trưng} & \textbf{Ý nghĩa} & \textbf{Công thức} \\
\midrule
Trung bình mẫu (\(\bar{x}\)) & Giá trị đại diện cho mẫu & \(\displaystyle \bar{x} = \frac{1}{n}\sum_{i=1}^n x_i\) \\
Phương sai mẫu (\(s^2\)) & Độ phân tán của dữ liệu & \(\displaystyle s^2 = \frac{1}{n-1}\sum (x_i - \bar{x})^2\) \\
Độ lệch chuẩn mẫu (\(s\)) & Căn bậc hai của phương sai & \(s = \sqrt{s^2}\) \\
Tỷ lệ mẫu (\(\hat{p}\)) & Tỷ lệ phần tử có tính chất A & \(\displaystyle \hat{p} = \frac{\text{số phần tử có A}}{n}\) \\
\bottomrule
\end{tabular}
\end{table}

\begin{figure}[H]
\centering
\begin{tikzpicture}
    \begin{axis}[
        width=10cm, height=6cm,
        title={Ví dụ: Phân bố chiều cao của mẫu 500 người},
        xlabel={Chiều cao (cm)},
        ylabel={Tần số},
        xmin=140, xmax=190,
        ymin=0, ymax=80,
        xtick={145,150,155,160,165,170,175,180,185},
        ytick=\empty,
        grid=both
    ]
    \addplot[ybar, fill=blue!30, bar width=3] coordinates {
        (147,5) (152,15) (157,35) (162,55) (167,70) (172,65) (177,40) (182,20) (187,8)
    };
    \draw[dashed, red, thick] (axis cs:167,0) -- (axis cs:167,75);
    \node[red] at (axis cs:167,78) {\(\bar{x} \approx 167\) cm};
    \end{axis}
\end{tikzpicture}
\caption{Biểu đồ phân bố chiều cao của mẫu, với trung bình mẫu \(\bar{x}\).}
\end{figure}

\begin{notebox}
\textbf{Lưu ý quan trọng:} Các đặc trưng mô tả chỉ có ý nghĩa trên \textbf{mẫu} mà bạn đã thu thập. Chúng chưa cho bạn biết điều gì về \textbf{tổng thể} lớn hơn. Để làm được điều đó, ta cần đến \textbf{thống kê suy luận}.
\end{notebox}

% =========================================================
% PHẦN 3: THỐNG KÊ SUY LUẬN
% =========================================================
\section*{\color{maincolor}III. Thống kê suy luận – Từ mẫu đến tổng thể}

\begin{defn}
\textbf{Thống kê suy luận} là tập hợp các phương pháp dùng để \textbf{suy luận} về các tham số của \textbf{tổng thể} dựa trên dữ liệu thu thập được từ \textbf{mẫu}.
\end{defn}

\begin{metabox}
Tiếp tục câu chuyện giám đốc kinh doanh: Bạn đã khảo sát 500 khách hàng và thấy \textbf{85\% hài lòng}. Câu hỏi đặt ra:

\vspace{0.3cm}
\textit{“Liệu có thể kết luận rằng toàn bộ hàng triệu khách hàng đều có tỷ lệ hài lòng khoảng 85\% không?”}

\vspace{0.3cm}
\textbf{Thống kê suy luận} sẽ giúp bạn trả lời câu hỏi này bằng cách:
\begin{itemize}
    \item Xây dựng \textbf{khoảng tin cậy} (confidence interval) cho tỷ lệ hài lòng của tổng thể.
    \item Thực hiện \textbf{kiểm định giả thuyết} (hypothesis testing) để xem có bằng chứng thống kê cho kết luận hay không.
\end{itemize}
\end{metabox}

\begin{figure}[H]
\centering
\begin{tikzpicture}[scale=0.8]
    \draw[thick, fill=green!20] (0,0) circle (3);
    \draw[thick, fill=blue!30] (0,0) circle (1.5);
    \node at (0,0) {Mẫu\\\(n = 500\)};
    \node at (0,3.5) {Tổng thể\\\(N = 100\) triệu};
    
    \draw[-{Stealth[length=3mm]}, thick, red] (1.5,0) -- (2.8,0);
    \node[red] at (3.5,0.5) {Suy luận};
    
    \draw[-{Stealth[length=3mm]}, thick, blue] (-1.5,0) -- (-2.8,0);
    \node[blue] at (-3.5,0.5) {Lấy mẫu};
\end{tikzpicture}
\caption{Mối quan hệ giữa tổng thể và mẫu trong thống kê suy luận.}
\end{figure}

\subsection*{3.1 Ví dụ kinh điển: Ước lượng chiều cao trung bình của dân số Việt Nam}

Giả sử chúng ta muốn biết \textbf{chiều cao trung bình của 100 triệu dân Việt Nam}. Không thể đo tất cả, chúng ta:

\begin{enumerate}
    \item Chọn ngẫu nhiên \(n = 500\) người (mẫu).
    \item Đo chiều cao của họ, tính \(\bar{x} = 167\) cm, \(s = 8\) cm.
    \item Dùng thống kê suy luận để ước lượng chiều cao trung bình của \textbf{cả nước}.
\end{enumerate}

\begin{theorem}
\textbf{Khoảng tin cậy 95\% cho trung bình tổng thể \(\mu\):}
\[
\boxed{\bar{x} \pm z_{\alpha/2} \cdot \frac{s}{\sqrt{n}}}
\]
Với \(n = 500\), \(\bar{x} = 167\), \(s = 8\), ta có:
\[
167 \pm 1.96 \times \frac{8}{\sqrt{500}} \approx 167 \pm 0.7 \Rightarrow (166.3, 167.7)
\]
\end{theorem}

\begin{bizbox}
\textbf{Kết luận kinh doanh:} Với độ tin cậy 95\%, chiều cao trung bình của người Việt Nam nằm trong khoảng từ \textbf{166.3 cm đến 167.7 cm}. 

Nếu bạn là một công ty sản xuất ghế, bạn sẽ thiết kế ghế phù hợp với khoảng chiều cao này!
\end{bizbox}

\begin{figure}[H]
\centering
\begin{tikzpicture}
    \begin{axis}[
        width=10cm, height=5cm,
        axis lines=middle,
        xlabel={Chiều cao (cm)},
        ylabel={Mật độ},
        domain=155:175,
        samples=100,
        title={Phân phối mẫu của \(\bar{x}\) và khoảng tin cậy 95\%}
    ]
    \addplot[thick, blue] {exp(-(x-167)^2/(2*(8/sqrt(500))^2))/(sqrt(2*pi)*(8/sqrt(500)))};
    
    \draw[dashed, red] (166.3,0) -- (166.3,0.5);
    \draw[dashed, red] (167.7,0) -- (167.7,0.5);
    \draw[<->, red, thick] (166.3,0.45) -- (167.7,0.45) node[midway, above] {KTC 95\%};
    
    \addplot[fill=red!20, domain=166.3:167.7] {exp(-(x-167)^2/(2*(8/sqrt(500))^2))/(sqrt(2*pi)*(8/sqrt(500)))} \closedcycle;
    \end{axis}
\end{tikzpicture}
\caption{Khoảng tin cậy 95\% cho chiều cao trung bình của tổng thể.}
\end{figure}

% =========================================================
% PHẦN 4: BẢN CHẤT CỦA VIỆC HỌC THỐNG KÊ
% =========================================================
\section*{\color{maincolor}IV. Học thống kê để làm gì? – Đọc tham số và ra quyết định}

\begin{center}
\begin{tikzpicture}
    \draw[thick, fill=maincolor!10] (0,0) rectangle (12,2.5);
    \node[draw, rounded corners, fill=maincolor!20, text width=2.5cm, align=center] at (2,1.2) {Thu thập\\dữ liệu};
    \node[draw, rounded corners, fill=maincolor!25, text width=2.5cm, align=center] at (5,1.2) {Mô tả\\mẫu};
    \node[draw, rounded corners, fill=maincolor!30, text width=2.5cm, align=center] at (8,1.2) {Suy luận\\tổng thể};
    \node[draw, rounded corners, fill=green!20, text width=2.5cm, align=center] at (11,1.2) {Quyết\\định};
    
    \draw[-{Stealth[length=3mm]}, thick] (3.3,1.2) -- (4.3,1.2);
    \draw[-{Stealth[length=3mm]}, thick] (6.3,1.2) -- (7.3,1.2);
    \draw[-{Stealth[length=3mm]}, thick] (9.3,1.2) -- (10.3,1.2);
\end{tikzpicture}
\captionof{figure}{Quy trình ra quyết định dựa trên thống kê.}
\end{center}

\begin{aibox}
\textbf{Vai trò của AI trong quy trình này:}

\begin{itemize}
    \item AI sẽ \textbf{tính toán} tất cả các tham số thống kê: \(\bar{x}\), \(s^2\), khoảng tin cậy, p-value...
    \item AI sẽ \textbf{vẽ biểu đồ}, \textbf{trực quan hóa} dữ liệu.
    \item AI sẽ \textbf{giải thích} ý nghĩa của các con số bằng ngôn ngữ dễ hiểu, kèm ẩn dụ.
    \item \textbf{Nhiệm vụ của bạn:} Đọc hiểu các tham số đó và \textbf{ra quyết định kinh doanh}.
\end{itemize}
\end{aibox}

\begin{bizbox}
\textbf{Ví dụ thực tế: Ra quyết định sản xuất}

Một công ty muốn tung ra sản phẩm mới. Họ khảo sát 1000 khách hàng và thấy 62\% “có khả năng mua”. AI tính khoảng tin cậy 95\% cho tỷ lệ này: (59\%, 65\%).

\begin{itemize}
    \item \textbf{Nếu} cận dưới 59\% vẫn cao hơn ngưỡng thành công 50\% → \textbf{quyết định sản xuất}.
    \item \textbf{Nếu} cận trên 65\% vẫn thấp hơn ngưỡng 70\% → \textbf{quyết định không sản xuất} hoặc cải tiến sản phẩm.
\end{itemize}

\textbf{AI không thể thay bạn ra quyết định. AI chỉ cho bạn con số. Bạn là người quyết định.}
\end{bizbox}

% =========================================================
% PHẦN 5: CÁCH VIẾT GIÁO TRÌNH MỚI
% =========================================================
\section*{\color{maincolor}V. Viết giáo trình thống kê theo phương pháp AI-first}

Từ cuộc hội thoại, chúng ta rút ra các nguyên tắc để xây dựng một giáo trình thống kê hoàn toàn mới:

\begin{enumerate}
    \item \textbf{AI làm thay mọi tính toán:} Sinh viên không cần nhớ công thức, không cần tính tay.
    \item \textbf{Tập trung vào bản chất:} Giải thích \textit{tại sao} có công thức đó, \textit{ý nghĩa} của từng tham số.
    \item \textbf{Ẩn dụ phong phú:} Dùng hình ảnh, câu chuyện để người học dễ hình dung.
    \item \textbf{Gắn với quyết định thực tế:} Mỗi bài học kết thúc bằng một tình huống kinh doanh cần ra quyết định.
    \item \textbf{Sinh viên tự thiết kế lộ trình:} Dùng AI để soạn bài, soạn slide, tự đặt câu hỏi dưới sự hướng dẫn của thầy.
\end{enumerate}

\begin{figure}[H]
\centering
\begin{tikzpicture}[node distance=0.6cm]
    \node[draw, rounded corners, fill=blue!10, text width=4cm, align=center] (step1) {Bước 1:\\Đưa đề cương + bài tập\\cho AI};
    \node[draw, rounded corners, fill=blue!15, text width=4cm, align=center, right=of step1] (step2) {Bước 2:\\AI viết bài giảng\\theo tinh thần bản chất};
    \node[draw, rounded corners, fill=blue!20, text width=4cm, align=center, below=of step2] (step3) {Bước 3:\\Sinh viên đọc,\\hiểu, đặt câu hỏi};
    \node[draw, rounded corners, fill=green!20, text width=4cm, align=center, left=of step3] (step4) {Bước 4:\\Ra quyết định\\trong tình huống thực};
    
    \draw[-{Stealth[length=3mm]}, thick] (step1) -- (step2);
    \draw[-{Stealth[length=3mm]}, thick] (step2) -- (step3);
    \draw[-{Stealth[length=3mm]}, thick] (step3) -- (step4);
    \draw[-{Stealth[length=3mm]}, thick] (step4) -- (step1);
\end{tikzpicture}
\caption{Vòng lặp học tập theo phương pháp AI-first.}
\end{figure}

% =========================================================
% PHẦN 6: TỔNG KẾT
% =========================================================
\section*{\color{maincolor}VI. Tổng kết bài học đầu tiên}

\begin{notebox}
\textbf{Những điều cần ghi nhớ:}

\begin{enumerate}
    \item \textbf{Thống kê có hai phần:} Thống kê mô tả (làm việc trên mẫu) và Thống kê suy luận (suy luận ra tổng thể).
    \item \textbf{Mẫu \(\neq\) Tổng thể:} Đặc trưng mẫu chỉ là ước lượng cho tham số tổng thể, kèm theo sai số.
    \item \textbf{Khoảng tin cậy} cho biết một phạm vi giá trị có khả năng chứa tham số tổng thể.
    \item \textbf{AI làm tính toán, con người ra quyết định.}
    \item \textbf{Người thầy thiết kế lộ trình, AI thực thi, sinh viên thực hành ra quyết định.}
\end{enumerate}
\end{notebox}

\begin{center}
\begin{tikzpicture}
    \draw[thick, fill=maincolor!5] (0,0) rectangle (12,2);
    \node[draw, rounded corners, fill=maincolor!20, text width=3cm, align=center] at (2,1) {Tổng thể\\\(N\) lớn};
    \node[draw, rounded corners, fill=maincolor!25, text width=3cm, align=center] at (6,1) {Mẫu\\\(n\) nhỏ};
    \node[draw, rounded corners, fill=green!20, text width=3cm, align=center] at (10,1) {Quyết định\\kinh doanh};
    
    \draw[-{Stealth[length=3mm]}, thick] (3.5,1) -- (4.5,1);
    \draw[-{Stealth[length=3mm]}, thick] (7.5,1) -- (8.5,1);
\end{tikzpicture}
\captionof{figure}{Bản chất của thống kê ứng dụng: Từ tổng thể → mẫu → quyết định.}
\end{center}

\begin{aibox}
\textbf{Lời nhắn từ Giáo sư AI:}

Các em đừng sợ thống kê. Các em đừng mất thời gian nhớ công thức. Hãy để AI làm việc đó. Việc của các em là \textbf{hiểu} và \textbf{quyết định}. 

Một người làm kinh doanh giỏi không phải người tính nhanh nhất, mà là người \textbf{đọc hiểu số liệu đúng} và \textbf{ra quyết định đúng lúc}. 

Hãy học thống kê như một giám đốc, không phải như một chiếc máy tính.
\end{aibox}

\begin{center}
    {\large\color{maincolor}\textbf{--- HẾT BÀI 1 ---}}
\end{center}

\begin{center}
    \textit{Bài tập về nhà: Hãy dùng AI để tính trung bình, phương sai và khoảng tin cậy 95\% cho bộ dữ liệu chiều cao của 10 người trong gia đình hoặc nhóm bạn. Sau đó, viết một đoạn văn ngắn (3-5 câu) đưa ra quyết định dựa trên kết quả đó.}
\end{center}

\vspace{1cm}

% =========================================================
% CÔNG CỤ TƯƠNG TÁC TRỰC TUYẾN
% =========================================================
\begin{tcolorbox}[
    colback=blue!3,
    colframe=maincolor,
    title=\faLaptop\ Công cụ thực hành trực tuyến,
    fonttitle=\bfseries,
    rounded corners,
    boxrule=1pt
]
\textbf{Hỗ trợ Quyết định Quản trị Nhân sự} -- Ứng dụng tương tác minh họa cách sử dụng thống kê mô tả và suy luận để ra quyết định kinh doanh thực tế.

\vspace{0.3cm}
Truy cập công cụ tại:

\begin{center}
    \Large\textbf{\href{https://nkoxteo123.github.io/MAS301/hr_decision_support.html}{\textcolor{maincolor}{nkoxteo123.github.io/MAS301/hr\_decision\_support.html}}}
\end{center}

\vspace{0.3cm}
\textbf{Hướng dẫn sử dụng:}
\begin{enumerate}
    \item Điều chỉnh các thanh trượt: trung bình mẫu (\(\bar{x}\)), độ lệch chuẩn (\(s\)), cỡ mẫu (\(n\)) và ngưỡng kỳ vọng (\(L\)).
    \item Quan sát biểu đồ hình chuông thay đổi theo thời gian thực.
    \item Đọc phần \textbf{giải thích logic} để hiểu tại sao hệ thống đưa ra quyết định đó.
    \item Thử các kịch bản khác nhau: Yếu đều, Phân hóa, Ổn định, Rủi ro tiềm ẩn.
\end{enumerate}
\end{tcolorbox}

\end{document}
